\documentclass[a4paper]{article}
\usepackage{graphicx}
\usepackage{amsmath,amssymb}
\usepackage{minted}
\usepackage{multirow}
\usepackage{float}
\usepackage{algorithm}
\usepackage{algpseudocode}
\usepackage{todonotes}
\usepackage{forest}
\usepackage{xstring}
\usepackage{xcolor}
\usepackage[most]{tcolorbox}
\usepackage{xparse}
\usepackage{natbib}
\usetikzlibrary{graphs,graphdrawing}
\usegdlibrary{layered}
% ND Mind Mirror default input files.
\input{/home/donkarlo/Dropbox/repo/nd_language_ai_project/src/nd_language_ai/formal/computer/programming/tex/latex/code_clips/external_graphviz.tex}
\input{/home/donkarlo/Dropbox/repo/nd_language_ai_project/src/nd_language_ai/formal/computer/programming/tex/latex/part/control_sequence/kind/semtag/semtag.tex}
\input{/home/donkarlo/Dropbox/repo/nd_language_ai_project/src/nd_language_ai/formal/computer/programming/tex/latex/code_clips/seehere.tex}

\title{Change-Point Detection}
\date{}
\author{Mohammad Rahmani}
\begin{document}
\maketitle

\section{Definition}
Change-point detection seeks times where the distribution or structure of a time series changes. A structural break can mean a different mean, variance, trend slope, or other model parameters. A one-day anomaly is not necessarily a change point.

\section{Segmentation Objective}
For boundaries $0=\tau_0<\tau_1<\cdots<\tau_K<\tau_{K+1}=n$, a widely used penalized objective is
\[
 \underset{K,\tau_1,\ldots,\tau_K}{\operatorname{minimize}}
 \left\{\sum_{k=0}^{K}
 C\bigl(y_{\tau_k+1:\tau_{k+1}}\bigr)
 +\lambda K\right\}.
\]
Here $C$ is a segment cost, such as squared errors around a segment mean; $\lambda$ penalizes adding too many breakpoints. To detect slope changes rather than level shifts, choose a cost based on fitted trends.

\section{Representative Algorithms}
\begin{itemize}
 \item \textbf{CUSUM:} cumulative-sum monitoring for persistent shifts, often used sequentially.
 \item \textbf{Binary segmentation:} recursively splits the sequence at candidate change points; efficient but approximate.
 \item \textbf{PELT:} exact penalized segmentation by dynamic programming with pruning under suitable assumptions on the cost. Linear expected time can occur in favorable conditions, but runtime is not universally linear.
\end{itemize}

\section{E-Commerce Example}
A redesign of a shopping-cart page may lead to a lasting conversion-rate shift. A detected breakpoint can prompt investigation of that release. A short holiday sales surge, by contrast, may reflect seasonality rather than a permanent regime change.

\section{Python Illustration}
\begin{verbatim}
import ruptures as rpt

# signal: one-dimensional ordered observations.
model = rpt.Pelt(model="l2", min_size=7).fit(signal)
end_indices = model.predict(pen=10)
# Typically the final end index is included as a
# boundary marker, not a discovered change point.
\end{verbatim}

\section{Advantages}
\begin{itemize}
 \item Detects changes missed by a single global regression slope.
 \item Supports operational investigations and monitoring.
 \item Different costs target different change types.
\end{itemize}

\section{Limitations}
\begin{itemize}
 \item Sensitive to penalty, minimum segment length, and the cost model.
 \item Seasonality and autocorrelation may generate false breaks.
 \item A detected breakpoint does not prove what caused the change.
 \item Offline segmentation can use future observations unavailable to live detectors.
\end{itemize}

\section{Practical Guidance}
Account for weekly effects, promotions, and outages before fitting. Compare detected dates with known business interventions and test sensitivity to penalty choices.

\bibliographystyle{plainnat}
\bibliography{/home/donkarlo/Dropbox/repo/nd_shared_research/bibliography/bibliography.bib}
\end{document}
